\newpage\thispagestyle{hl}
\hlbanner{Part II $\cdot$ Full-order physics}{6}{Computational acceleration: vectorised assembly and batched kernels}
\hlsection{The question}
\begin{hlquote}Profile a nonlinear run on a plate of a couple of thousand elements and nearly all the time goes into assembly, and nearly all of that into the Python interpreter, not into arithmetic.\end{hlquote}
\hlsection{Key concepts}
\begin{description}\setlength\itemsep{5pt}\setlength\parsep{0pt}\renewcommand{\descriptionlabel}[1]{\hspace{\labelsep}\sffamily\bfseries\color{hlnavy}#1}
\item[Vectorisation] array operations in compiled code replace the loop over elements
\item[Batched kernel] all elements of a block are computed in one call
\item[Sparse tangent] the tangent stays sparse so Newton can use a sparse factorisation
\item[Chunking] the batch size is set to bound memory
\item[Support report] says which blocks cannot be batched, and why
\item[PyTorch backend] optional; not run where the book is built, and said so
\end{description}
\hlsection{After this chapter you can}
\begin{itemize}\setlength\itemsep{3pt}
\item decide when the batched path pays off and check it against the loop
\item read the support report and set the chunk size to bound memory
\item say what the optional PyTorch backend replaces, and at which sizes it pays off
\end{itemize}
\hlsection{In this chapter}
{\begin{multicols}{2}\raggedright
\textbf{6.1}\ Where the time goes\\
\textbf{6.2}\ What the batched path does\\
\textbf{6.3}\ Using it\\
\textbf{6.4}\ Verification\\
\textbf{6.5}\ The optional PyTorch backend\\
\textbf{6.6}\ Related pieces\\
\textbf{6.7}\ Limits\\
\textbf{6.8}\ Summary
\end{multicols}}
\hlsection{Explained on the page, in short boxes}
\noindent Interpreted loops and vectorisation; Assembly as a scatter-add; Batching and chunking; GPUs and floating-point precision.
\newpage\thispagestyle{hl}
\noindent{\sffamily\small\color{hlgrey}Chapter 6 \quad Computational acceleration: vectorised assembly and batched kernels}\par
\hlfig{\hlimg{ch06_scaling}}{Assembly time of the element loop and the batched route against the number of elements (left), and the ratio of the two (right).}
\hlsection{Headline result}
\begin{hlheadline}A whole Newton solve on a plate of \cbkNewtonNE{} elements and \cbkNewtonN{} unknowns takes \cbkNewtonBatched{} s on the batched path against \cbkNewtonDefault{} s on the default path, a factor of \cbkNewtonRatio{}, and the solutions agree to \cbkNewtonDiff.\end{hlheadline}
\hlsection{Main classes and functions}
{\begin{itemize}\setlength\itemsep{2pt}
\item \texttt{fea\_engine.solver.FESystem}
\item \texttt{fea\_engine.newton\_options.NewtonOptions}
\item \texttt{fea\_engine.elements.nonlinear\_solids.Tet10SolidTL}
\item \texttt{fea\_engine.torch\_sparse\_solver.torch\_sparse\_cg\_solve}
\end{itemize}}
\hlsection{A stated limit}
\begin{hllimit}The gain depends on the element: arithmetic-heavy hexahedra gain less than interpreter-bound plane elements, and once assembly is fast the sparse solve can dominate.\end{hllimit}
\newpage\thispagestyle{hl}
\noindent{\sffamily\small\color{hlgrey}Chapter 6 \quad Computational acceleration: vectorised assembly and batched kernels}\par
\hlsection{Highlight: the optional PyTorch patch}
\begin{hlquote}Nothing in the default path changes. A call without the two extra arguments runs the old code, and PyTorch is imported only when a torch route is asked for.\end{hlquote}
\hlfig{\hltikz{ch06_torch}}{How the PyTorch backend is patched in. A call with \texttt{backend="torch"} leaves the default NumPy and SciPy route untouched and takes a separate path of three stages: conversion to tensors on the chosen device, computation in double precision, and conversion back to NumPy.}
\hlsection{Three routes behind one switch}
\begin{itemize}\setlength\itemsep{3pt}
\item \textbf{Dense solve:} \texttt{torch.linalg.solve} on a dense tensor; any matrix, small systems and an independent reference.
\item \textbf{Sparse conjugate gradients:} Jacobi-preconditioned CG on a sparse CSR tensor; needs a symmetric positive definite matrix, and a non-positive diagonal is refused.
\item \textbf{Stress recovery and element kernels:} written once for an array namespace \texttt{xp}, so the NumPy and torch versions differ only in which module is passed.
\end{itemize}
\hlsection{What it leaves alone, and what it does not claim}
\begin{hllimit}Assembly of $\bm K$ and $\bm F$ and the Newton drivers stay in NumPy, so a nonlinear analysis does not run on the GPU through this backend. PyTorch is not installed where the book is built: the CPU and GPU timings are records of the author's machine, not measured for the book. What was run here is the algebra: the kernels called with \texttt{xp=numpy} reproduce the batched element kernels to \bktShadowForce{} (force) and \bktShadowTangent{} (tangent).\end{hllimit}

